% ============================================================
% s02b-chI-10.tex —— Chapter I · §I.10「谱作用」正文部分（p.37--p.38）
% 底本：arXiv:2102.13459v4，PDF p.37--p.38
% 说明：§I.10 起于 p.37（本文件），p.39 起是 Remark I.10.3，由块 C 的
%       s03-chI-3.tex 承接；本文件自开 \section，不要在两处重复开节。
% 分工缘由：块 B 的任务边界定在 p.36（§I.9 末），块 C 从 p.39 起，
%           p.37--38 由 Lead 补译（2026-11-23）。
% ============================================================

\section{谱作用}

\srcpage{37}
我们构造的范畴结构其实给出了更好的东西。设 $\Lambda$ 是 $\QQ_\ell(\sqrt q)$ 的某个有限扩张的整数环。
我们有紧对象的稳定 $\infty$-范畴 $\mathcal C=D_{\mathrm{lis}}(\Bun_G,\Lambda)^{\omega}$，它在 $\Lambda$ 上是线性的；
并且对有限集 $I$ 函子性地有一个正合幺半函子
\[
  \Rep_\Lambda(\hat G\rtimes Q)^I\longrightarrow \operatorname{End}_\Lambda(\mathcal C)^{BW_E^I},
\]
它关于 $\Rep_\Lambda(Q^I)$ 是线性的。这里 $\operatorname{End}_\Lambda(\mathcal C)$ 表示 $\mathcal C$ 的 $\Lambda$-线性自函子所成的稳定 $\infty$-范畴；我们借助把 $\mathcal C$ 看作在相对离散的凝聚态 $\Lambda$-模中丰富，而把它看作在凝聚态 $\Lambda$-模中丰富。\yizhu{「相对离散的凝聚态 $\Lambda$-模」即把抽象 $\Lambda$-模按离散拓扑视为凝聚态模；原文此处只为说明 enrichment 的来源，不影响后文定理的陈述。}{A：原文 §I.10 p.37；glossary.tsv（condensed＝凝聚态）}
下一定理的第一个版本属于 Nadler--Yun [NY19]（Betti 几何 Langlands 的情形），更一般的版本出现在 Gaitsgory--Kazhdan--Rozenblyum--Varshavsky [GKRV22] 的工作中。
然而这两处文献实际上都假设 $G$ 分裂、只处理特征 $0$ 系数，并以一个离散群代替 $W_E$。
至少向 $\ZZ_\ell$-系数作推广是一件非平凡的事。\yizhu{上标 $BW_E^I$ 在引言中未给出定义。按第 IX 章（Hecke 作用）的用法，它标出的是自函子范畴上 $W_E^I$ 的 Bernstein--Whittaker 型等变结构；此处照原文保留记号，读者可对照第 IX 章。}{I：由 Ch IX §IX.2 的记号与本节上下文推断，原文本节未定义}

注意 $Z^1(W_E,\hat G)$ 不是拟紧的，因为它有无穷多个连通分支；它可以写成开闭拟紧子概形 $Z^1(W_E/P,\hat G)$ 的递增并。
我们说 $\operatorname{Perf}(Z^1(W_E,\hat G)/\hat G)$ 在稳定 $\infty$-范畴 $\mathcal C$ 上的一个作用具有\textbf{紧支撑}，如果对所有 $X\in\mathcal C$，函子
$\operatorname{Perf}(Z^1(W_E,\hat G)/\hat G)\to\mathcal C$（由作用在 $X$ 上诱导）经由某个 $\operatorname{Perf}(Z^1(W_E/P,\hat G)/\hat G)$ 分解。\yizhu{「紧支撑」在这里是类比性的说法：$Z^1(W_E,\hat G)$ 非拟紧，只能沿 $Z^1(W_E/P,\hat G)$ 的递增并逐块作用，故要求每个对象的像只落在有限一段上，恰似支撑集紧。}{I：由上文「递增并」与定义本身的量化结构推出}

\begin{theorem}% 原编号 I.10.1
假设 $\ell$ 不整除 $\pi_0Z(G)$ 的阶。设 $\mathcal C$ 是小 $\Lambda$-线性稳定 $\infty$-范畴。
则「对有限集 $I$ 函子性地给出一个正合 $\Rep_\Lambda(Q^I)$-线性幺半函子
\[
  \Rep_\Lambda(\hat G\rtimes Q)^I\longrightarrow \operatorname{End}_\Lambda(\mathcal C)^{BW_E^I}
\]
」等价于「给出
\[
  \operatorname{Perf}(Z^1(W_E,\hat G)_\Lambda/\hat G)
\]
的一个紧支撑 $\Lambda$-线性作用」。

这里，给定 $\operatorname{Perf}(Z^1(W_E,\hat G)_\Lambda/\hat G)$ 的一个紧支撑 $\Lambda$-线性作用，可以对 $I$ 函子性地造出这样的正合 $\Rep_\Lambda(Q^I)$-线性幺半函子
\[
  \Rep_\Lambda(\hat G\rtimes Q)^I\longrightarrow \operatorname{End}_\Lambda(\mathcal C)^{BW_E^I}：
\]
办法是把正合 $\Rep_\Lambda(Q^I)$-线性对称幺半函子
\[
  \Rep_\Lambda(\hat G\rtimes Q)^I\longrightarrow \operatorname{Perf}(Z^1(W_E,\hat G)_\Lambda/\hat G)^{BW_E^I}
\]
与 $\operatorname{Perf}(Z^1(W_E,\hat G)_\Lambda/\hat G)$ 的作用复合。\yizhu{「对 $I$ 函子性」是这段的技术要点：$I$ 只出现在对称幺半函子的源与靶上，作用本身与 $I$ 无关，故作用一个就得到整族函子。}{I：由陈述中 $I$ 只出现于两侧上标与源范畴看出}

此结论对 $\Lambda$ 取 $\QQ_\ell$ 上的域、任意素数 $\ell$ 同样成立。
\end{theorem}

这里，正合 $\Rep_\Lambda(Q^I)$-线性对称幺半函子
\[
  \Rep_\Lambda(\hat G\rtimes Q)^I\longrightarrow \operatorname{Perf}(Z^1(W_E,\hat G)_\Lambda/\hat G)^{BW_E^I}
\]
由张量积以及下面的正合 $\Rep_\Lambda(Q)$-线性对称幺半函子诱导：
\[
  \Rep_\Lambda(\hat G\rtimes Q)\longrightarrow \operatorname{Perf}(Z^1(W_E,\hat G)_\Lambda/\hat G)^{BW_E},
\]
后者对应于万有 $\hat G\rtimes Q$-主丛，其中万有的 $W_E$-等变性由 $Z^1(W_E,\hat G)/\hat G$ 参数化。

\srcpageinl{38}证明的关键部分其实是上面定理 I.8.2 的最后一部分，它实际上以生成元与关系给出了 $\operatorname{Perf}(Z^1(W_E/P,\hat G)/\hat G)$ 的描述，本定理也是如此。
特别地，我们得到 $\operatorname{Perf}(Z^1(W_E,\hat G)_\Lambda/\hat G)$ 在 $D_{\mathrm{lis}}(\Bun_G,\Lambda)$ 上的作用，它与 Hecke 作用适度相容。

万事俱备，现在显然主要猜想如下，参看 [AG15]、[BZCHN20]、[Zhu20]、[Hel23]：\footnote{本手稿的上一版本提出的猜想看起来不够精确：它只要求存在一个函子，而没有注意到该函子必定实现为右伴随、因而若存在则唯一。这一唯一性尤其由 Hansen [Han24] 向我们指出。}

\begin{conjecture}% 原编号 I.10.2
假设 $G$ 拟分裂，并选取 Whittaker 数据：一个 Borel 子群 $B\subset G$，以及单幂根 $U\subset B$ 的一个一般特征标 $\psi:U(E)\to\cO_L^\times$，其中 $L/\QQ_\ell$ 是某个代数扩张；再固定 $\sqrt q\in\cO_L$。
设 $n$ 是 $\pi_0Z(G)$ 的阶，$\Lambda=\cO_L[\tfrac1n]$。设
\[
  \mathcal{W}_\psi\in D_{\mathrm{lis}}(\Bun_G,\Lambda)
\]
是 Whittaker 层，即集中在 $\Bun_G^1$ 上、对应于 Whittaker 表示 $c\text{-}\mathrm{Ind}_{U(E)}^{G(E)}\psi$ 的层；并设
\[
  \operatorname{Ind}\operatorname{Perf}^{\mathrm{qc}}(Z^1(W_E,\hat G)_\Lambda/\hat G)\longrightarrow D_{\mathrm{lis}}(\Bun_G,\Lambda):\quad M\longmapsto \mathrm{Act}_M(\mathcal{W}_\psi)
\]
定义为谱作用 $\mathrm{Act}$ 在 $\mathcal{W}_\psi$ 上的保余极限延拓。
则相应的右伴随函子在限制到紧对象时全忠实，并诱导 $(\operatorname{Perf}(Z^1(W_E,\hat G)_\Lambda/\hat G)$-线性小稳定$)$ $\infty$-范畴的等价
\[
  \mathcal D(\Bun_G,\Lambda)^{\omega}\simeq \mathcal D^{b,\mathrm{qc}}_{\mathrm{coh},\mathrm{Nilp}}(Z^1(W_E,\hat G)_\Lambda/\hat G).
\]
\end{conjecture}

我们在这里把 $\pi_0Z(G)$ 的阶 $n$ 反转了，因为只有这时谱作用才被构造出来。
如果没有反转 $n$，我们并不确定该期待什么。（事实上，如果我们在坏素数处所用的「幂零奇异支撑」这一概念需要修改，我们也不会感到意外。）

这里我们使用了带幂零奇异支撑的凝聚层复形这一概念，见 [AG15]。
更确切地说，$\mathcal D^{b,\mathrm{qc}}_{\mathrm{coh},\mathrm{Nilp}}$ 是具拟紧支撑、凝聚上同调与幂零奇异支撑的有界复形所成的 $\infty$-范畴。
在特征 $0$ 系数下，或在 banal 素数 $\ell$ 处，幂零奇异支撑这一条件实际上是自动的。\yizhu{banal 素数：指 $\ell$ 不整除相关有限群的阶、使 $p$-进表示论「不脱离特征 $0$ 直觉」的素数；中文无统一译名，故保留原文。}{I：本库无既有译法；原文此处未给定义}

若 $\mathcal{W}_\psi$ 是 Whittaker 层，并以 $*$ 记谱作用，则该猜想是说
\[
  \operatorname{Perf}^{\mathrm{qc}}(Z^1(W_E,\hat G)_\Lambda/\hat G)\longrightarrow \mathcal D(\Bun_G,\Lambda),\qquad
  M\longmapsto M*\mathcal{W}_\psi
\]
全忠实，并延拓为稳定 $\infty$-范畴的等价
\[
  \mathcal D^{b,\mathrm{qc}}_{\mathrm{coh},\mathrm{Nilp}}(Z^1(W_E,\hat G)_\Lambda/\hat G)\simeq \mathcal D(\Bun_G,\Lambda)^{\omega}.
\]

回忆右边全忠实地包含 $\mathcal D(G(E),\Lambda)^{\omega}$，所以特别地，这个 $\infty$-范畴应当全忠实地嵌入左边。
这一点已由 Hellmann 在 [Hel23] 中猜想，Ben-Zvi--Chen--Helm--Nadler [BZCHN20] 证明了一部分（他们使用 $\QQ_\ell$-系数，并处理分裂群以及与具有 Iwahori 不动向量的表示相对应的 Bernstein 分支）。
